\documentclass{article}
\usepackage[utf8]{inputenc}
\usepackage[portuguese]{babel}
\usepackage{graphicx}
\usepackage{geometry}
\usepackage{float}
\usepackage{amsmath}
\usepackage[dvipsnames]{xcolor}
\usepackage{listings}
\usepackage{caption}
\geometry{a4paper, margin=2.5cm}
\definecolor{codebg}{HTML}{eeeeee}
\definecolor{codeframe}{HTML}{000000}

\lstdefinestyle{pythoncode}{
    language=Python,
    basicstyle=\ttfamily,
    backgroundcolor=\color{codebg},
    frame=single,
    rulecolor=\color{codeframe},
    framerule=1pt,
    framesep=10pt,
    lineskip=0pt,
    breaklines=true,
    showstringspaces=false
}

\begin{document}

\begin{titlepage}
    \centering
    \vspace*{1cm}
    \includegraphics[width=0.4\textwidth]{fcup_logo.png}\par\vspace{1.5cm}

    {\Large \textbf{Faculdade de Ciências da Universidade do Porto}}\par
    \vspace{0.2cm}
    {\large Departamento de Matemática}\par
    \vspace{2cm}
    
    {\large \textbf{Métodos Numéricos (M2039)}}\par
    \vspace{0.5cm}
    {\large Ano Letivo 2026/2027}\par
    \vspace{2cm}
    
    {\huge \textbf{Trabalho Prático 1}}\par
    \vspace{0.5cm}
    {\Large Relatório}\par
    \vspace{3cm}
    
    \textbf{Realizado por:}\par
    \vspace{0.5cm}
    \begin{tabular}{ll}
        Pedro André Silva Pinho Gonçalves & up202405642 \\
        Filipe Marques Mano & up202506000 \\
        Martim Andrade Ribeiro & up202505495 \\
    \end{tabular}
    
    \vfill
    
    {\large Outubro 2026}\par
\end{titlepage}

\section{Motivação do trabalho}

Este trabalho prático tem como foco a resolução de problemas de 
aproximação de constantes matemáticas (como $e$ e $\pi^2$) em ambiente 
computacional, exigindo que os resultados cumpram critérios de paragem 
analíticos baseados na majoração do erro de truncatura. Paralelamente, 
procuramos investigar o limite de precisão da máquina (epsilon da máquina)
para ver como este pode influenciar os resultados obtidos.

\pagebreak

\section{Resolução da Questão 2}
\smallskip
\subsection{Majorar o erro}

No âmbito da resolução da questão, é pedido que seja calculado 
um valor aproximado da soma da série dada por:

\begin{equation}
    S = 18 \sum_{k=1}^{\infty} \frac{k!^2}{k^2 (2k)!}
\end{equation}
\medskip

Inicialmente, procuramos definir uma simplificação para o termo 
geral da série S, dado por: $a_k = \frac{k!^2}{k^2 (2k)!}$, de modo 
a obter uma nova expressão através da qual possamos realizar o 
cálculo aproximado. Para isso, analisamos o termo seguinte, ou seja, 
$a_{k+1} = \frac{(k+1)!^2}{(k+1)^2 (2(k+1))!}$, e realizamos a 
divisão entre os dois termos, obtendo:

\begin{align*}
    \frac{a_{k+1}}{a_k} &= \frac{(k+1)!^2}{(k+1)^2 (2(k+1))!} \cdot \frac{k^2 (2k)!}{k!^2} \\
    &= \frac{(k+1)^2 k!^2}{(k+1)^2 (2k+2)(2k+1)(2k)!} \cdot \frac{k^2 (2k)!}{k!^2} \\
    &= \frac{k^2}{(2k+1)(2k+2)} \\
    &= \frac{k^2}{4k^2 + 6k + 2}
\end{align*}
\medskip

Após esta simplificação, deduzimos que o termo geral da série S 
pode ser reescrito como:

\begin{equation}
    a_{k+1} = a_k \cdot \frac{k^2}{4k^2 + 6k + 2}
\end{equation}
\bigskip

De modo a majorar o termo geral da série, podemos observar que, 
como k é um número positivo, temos $4k^2 + 6k + 2 > 4k^2$, o que 
nos permite escrever a seguinte desigualdade:
\medskip

\begin{equation}
    \frac{k^2}{4k^2 + 6k + 2} < \frac{k^2}{4k^2} = \frac{1}{4}
\end{equation}
\bigskip

Assim, temos que:

\begin{equation}
    \forall k \frac{a_{k+1}}{a_k} < \frac{1}{4} < 1
\end{equation}
\bigskip

Logo, podemos usar a seguinte formula para majorar o erro:
\medskip

\begin{equation}
    \Delta S_k \leq \frac{a_{k+1}}{1-\frac{1}{4}} = a_{k+1} \times \frac{4}{3}
\end{equation}
\pagebreak

\subsection{Primeiro algoritmo}

Inicialmente, decidimos implementar o algoritmo desta maneira:

\begin{center}
\begin{minipage}{0.8\linewidth}
\begin{lstlisting}[style=pythoncode]
def S(eps):

    # Variaveis
    acc = 0.0
    ter = 9.0
    err = math.inf
    k = 1

    # Calcular e somar termos
    while (err > eps and k < 100):
        acc += ter
        ter *= (k ** 2) / ((2*k + 1) * (2*k + 2))
        err = (4/3) * ter
        k += 1
        
    # Retornar o numero de iteracoes,
    # o resultado e o erro absoluto
    return k - 1, acc, abs(math.pi**2 - acc)
\end{lstlisting}
\end{minipage}
\end{center}

Após executar os testes pedidos no exercicio, obtivemos a seguinte tabela:

\begin{table}[H]
    \centering
    \begin{tabular}{|c|c|l|l|}
        \hline
        \textbf{$\epsilon$} & \textbf{Iterações} & \textbf{Resultado} & \textbf{Erro Absoluto Efetivo} \\ \hline
        $10^{-8}$  & 13 & 9.869604398132752 & $2.9566056269914043 \times 10^-09$ \\ \hline
        $10^{-9}$  & 14 & 9.869604400421997 & $6.673612773511195  \times 10^-10$ \\ \hline
        $10^{-10}$ & 16 & 9.869604401054712 & $3.4646063795662485 \times 10^-11$ \\ \hline
        $10^{-11}$ & 17 & 9.869604401081402 & $7.956302283673722  \times 10^-12$ \\ \hline
        $10^{-12}$ & 19 & 9.869604401088933 & $4.2454928461665986 \times 10^-13$ \\ \hline
        $10^{-13}$ & 21 & 9.869604401089337 & $2.1316282072803006 \times 10^-14$ \\ \hline
        $10^{-14}$ & 22 & 9.869604401089354 & $3.552713678800501  \times 10^-15$ \\ \hline
        $10^{-15}$ & 24 & 9.86960440108936  & $1.7763568394002505 \times 10^-15$ \\ \hline
    \end{tabular}
    \vspace{0.2cm}
    \caption{Resultados da aproximação da série S ($\pi^2$)}
\end{table}

Algo que podemos notar olhando para a tabela é que a função falha
no ultimo teste. Era esperado um valor com erro menor do que $10^{-15}$,
no entanto foi calculado um valor de erro superior. Como $\pi^2$ está
entre 8 e 16, o seu expoente em representação virgula flutuante binária
é igual a 3, e como nós estamos a trabalhar com dupla precisão, temos
que o erro por arredondamento é de até
$2^{3-53} = 2^{-50} \approx 9 \times 10^{-16}$, o que é muito próximo
de $10^{-15}$. Isto significa que é muito dificil de obter um erro menor
do que $10^{-15}$, mas não é impossivel.
\pagebreak

\subsection{Segundo algoritmo}

Modificando o programa para, em vez de somar os termos à medida que os
calcula, realizar as somas de menor magnitude primeiro de modo a
minimizar os erros por arredondamento, obtemos o seguinte código
e tabela:

\begin{center}
\begin{minipage}{0.8\linewidth}
\begin{lstlisting}[style=pythoncode]
def S(eps):

    # Variaveis
    lst = []
    ter = 9.0
    err = math.inf
    k = 1
    res = 0.0

    # Calcular termos
    while (err > eps and k < 25):
        lst.append(ter)
        ter *= (k**2) / ((2*k + 1) * (2*k + 2))
        err = (4/3) * ter
        k += 1

    # Somar em ordem crescente de magnitude
    lst.sort(key = abs)
    while len(lst) > 1:
        res = lst[0] + lst[1]
        lst = lst[2:]
        bisect.insort(lst, res, key = abs)
    res = lst[0]

    # Retornar o numero de iteracoes,
    # o resultado e o erro absoluto
    return (k - 1, res, abs(math.pi**2 - res))
\end{lstlisting}
\end{minipage}
\end{center}

\begin{table}[H]
    \centering
    \begin{tabular}{|c|c|l|l|}
        \hline
        \textbf{$\epsilon$} & \textbf{Iterações} & \textbf{Resultado} & \textbf{Erro Absoluto Efetivo} \\ \hline
        $10^{-8}$  & 13 & 9.86960439813275  & $2.9566074033482437 \times 10^-09$ \\ \hline
        $10^{-9}$  & 14 & 9.869604400421997 & $6.673612773511195 \times 10^-10$  \\ \hline
        $10^{-10}$ & 16 & 9.86960440105471  & $3.4647840152501885 \times 10^-11$ \\ \hline
        $10^{-11}$ & 17 & 9.8696044010814   & $7.958078640513122 \times 10^-12$  \\ \hline
        $10^{-12}$ & 19 & 9.869604401088933 & $4.2454928461665986 \times 10^-13$ \\ \hline
        $10^{-13}$ & 21 & 9.869604401089335 & $2.3092638912203256 \times 10^-14$ \\ \hline
        $10^{-14}$ & 22 & 9.869604401089353 & $5.329070518200751 \times 10^-15$  \\ \hline
        $10^{-15}$ & 24 & 9.869604401089358 & $0.0$                              \\ \hline
    \end{tabular}
    \vspace{0.2cm}
    \caption{Resultados da aproximação da série S ($\pi^2$)}
\end{table}

Se olharmos de novo para o ultimo teste, vêmos que o erro absoluto agora
é marcado como $0.0$, isto significa que o resultado é exatamente igual
à representação de $\pi^2$ em precisão dupla, ou seja, o erro associado
é o erro por arredondamento: $2^{-50} < 9 \times 10^{-16} < 10^{-15}$ \\
Assim, o nosso programa satisfaz todos os critérios do exercicio.
\pagebreak

\section{Resolução da Questão 3}

\subsection{Métodos}
Para a aproximação da constante de Euler ($e \approx 2.71828...$), o 
guião sugere a utilização de duas abordagens matemáticas distintas, 
a Fórmula A e a Fórmula B, definidas respetivamente por:

\[ \text{Fórmula A: } \quad a_n = \left(1 + \frac{1}{n}\right)^n \]
\[ \text{Fórmula B: } \quad b_m = \sum_{k=0}^{m} \frac{1}{k!} \]

A implementação da Fórmula A consiste na tradução direta da expressão 
algébrica numa função simples:

\begin{center}
\begin{minipage}{0.8\linewidth}
\begin{lstlisting}[style=pythoncode]
def a_euler(n):
    return (1 + 1/n)**n>
\end{lstlisting}
\end{minipage}
\end{center}

Já na Fórmula B, o algoritmo baseia-se num ciclo \texttt{while (k <= m)} 
que vai acumulando os termos da série a cada iteração. Para otimizar o 
código, evitamos o cálculo repetitivo de fatoriais do zero a cada iteração 
implementando uma relação de recorrência. Sabendo que o termo geral da 
série é $a_k = \frac{1}{k!}$, o termo seguinte pode ser deduzido a partir 
do termo atual através da expressão:
\[ a_{k+1} = \frac{1}{(k+1)!} = \frac{1}{(k+1) \times k!} = a_k \times \frac{1}{k+1} \]

Desta forma, o programa calcula o próximo valor a adicionar à soma 
multiplicando simplesmente o termo anterior por $\frac{1}{k+1}$:

\begin{center}
\begin{minipage}{0.8\linewidth}
\begin{lstlisting}[style=pythoncode]
def b_euler(m):
    soma = 0.0
    termo = 1.0
    k = 0
    while (k <= m):
        soma += termo
        termo = termo * (1 / (k + 1))
        k += 1
    return soma
\end{lstlisting}
\end{minipage}
\end{center}
\break

\subsection{Discussão dos Resultados}

Pela análise dos resultados da fórmula A (Tabela \ref{tab:formula_a}), 
observa-se que, ao aumentar o valor de $n$, a sucessão começa a convergir 
para o valor da constante de Euler. No entanto, a partir de $n=10^{12}$, 
nota-se o início de uma divergência progressiva em relação ao valor de 
referência. Esta alteração deve-se à aproximação ao limite de precisão 
da máquina, o que provoca perdas severas de algarismos significativos. 
Ao elevar esta base, já com erros de arredondamento, a um expoente muito 
elevado, o pequeno erro inicial sofre uma propagação massiva que corrompe 
o resultado final. O caso extremo ocorre quando $n=10^{16}$, onde a parcela 
$1/n$ ($10^{-16}$) é inferior ao epsilon da máquina. Nesta iteração, 
o computador não consegue distinguir $1.0$ da soma $1.0 + 10^{-16}$, 
efetuando o arredondamento para $1.0$. Por fim, a operação realizada 
é $1.0^{10^{16}}$, o que resulta em $1.0$.
\bigskip

\begin{table}[H]
    \centering
    \begin{tabular}{|c|l|l|}
        \hline
        \textbf{$n$} & \textbf{Resultado} & \textbf{Erro Absoluto Efetivo} \\ \hline
        $10^{1}$  & 2.5937424601000023 & $1.245 \times 10^{-1}$ \\ \hline
        $10^{2}$  & 2.7048138294215285 & $1.347 \times 10^{-2}$ \\ \hline
        $10^{3}$  & 2.7169239322355936 & $1.358 \times 10^{-3}$ \\ \hline
        $10^{4}$  & 2.7181459268249255 & $1.359 \times 10^{-4}$ \\ \hline
        $10^{5}$  & 2.7182682371922975 & $1.359 \times 10^{-5}$ \\ \hline
        $10^{6}$  & 2.7182804690957534 & $1.359 \times 10^{-6}$ \\ \hline
        $10^{7}$  & 2.7182816941320818 & $1.343 \times 10^{-7}$ \\ \hline
        $10^{8}$  & 2.7182817983473577 & $3.011 \times 10^{-8}$ \\ \hline
        $10^{9}$  & 2.7182820520115603 & $2.236 \times 10^{-7}$ \\ \hline
        $10^{10}$ & 2.7182820532347876 & $2.248 \times 10^{-7}$ \\ \hline
        $10^{11}$ & 2.7182820533571100 & $2.249 \times 10^{-7}$ \\ \hline
        $10^{12}$ & 2.7185234960372378 & $2.417 \times 10^{-4}$ \\ \hline
        $10^{13}$ & 2.7161100340869010 & $2.172 \times 10^{-3}$ \\ \hline
        $10^{14}$ & 2.7161100340870230 & $2.172 \times 10^{-3}$ \\ \hline
        $10^{15}$ & 3.0350352065492620 & $3.168 \times 10^{-1}$ \\ \hline
        $10^{16}$ & 1.0000000000000000 & $1.718 \times 10^{0}$  \\ \hline
    \end{tabular}
    \vspace{0.2cm}
    \caption{Resultados da aproximação da constante de Euler utilizando a Fórmula A.}
    \label{tab:formula_a}
\end{table}
\pagebreak

Por outro lado, a fórmula B (Tabela \ref{tab:formula_b}) mostra um 
comportamento estável. Ao aumentar a quantidade de termos da série, 
os valores convergem de forma consistente para o valor real da constante, 
verificando-se uma estagnação a partir da iteração $m=17$. Esta estagnação 
deve-se ao limite de precisão do computador. Para $m=18$, o termo a 
adicionar à soma é $1/18! \approx 1.56 \times 10^{-16}$, que é inferior 
ao epsilon da máquina. Por conseguinte, os termos a partir da décima 
oitava iteração não agregam qualquer valor à soma acumulada, uma vez que 
estes excedem a capacidade de representação da mantissa do sistema, sendo 
ignorados durante o processo de arredondamento.

\begin{table}[H]
    \centering
    \begin{tabular}{|c|l|l|}
        \hline
        \textbf{$m$} & \textbf{Resultado} & \textbf{Erro Absoluto Efetivo} \\ \hline
        $1$  & 2.0000000000000000 & $7.183 \times 10^{-1}$ \\ \hline
        $2$  & 2.5000000000000000 & $2.183 \times 10^{-1}$ \\ \hline
        $3$  & 2.6666666666666665 & $5.162 \times 10^{-2}$ \\ \hline
        $4$  & 2.7083333333333330 & $9.948 \times 10^{-3}$ \\ \hline
        $5$  & 2.7166666666666663 & $1.615 \times 10^{-3}$ \\ \hline
        $6$  & 2.7180555555555554 & $2.263 \times 10^{-4}$ \\ \hline
        $7$  & 2.7182539682539684 & $2.786 \times 10^{-5}$ \\ \hline
        $8$  & 2.7182787698412700 & $3.059 \times 10^{-6}$ \\ \hline
        $9$  & 2.7182815255731922 & $3.029 \times 10^{-7}$ \\ \hline
        $10$ & 2.7182818011463845 & $2.731 \times 10^{-8}$ \\ \hline
        $11$ & 2.7182818261984930 & $2.261 \times 10^{-9}$ \\ \hline
        $12$ & 2.7182818282861687 & $1.729 \times 10^{-10}$ \\ \hline
        $13$ & 2.7182818284467594 & $1.229 \times 10^{-11}$ \\ \hline
        $14$ & 2.7182818284582300 & $8.149 \times 10^{-13}$ \\ \hline
        $15$ & 2.7182818284589950 & $5.018 \times 10^{-14}$ \\ \hline
        $16$ & 2.7182818284590430 & $2.220 \times 10^{-15}$ \\ \hline
        $17$ & 2.7182818284590455 & $4.441 \times 10^{-16}$ \\ \hline
        $18$ & 2.7182818284590455 & $4.441 \times 10^{-16}$ \\ \hline
        $19$ & 2.7182818284590455 & $4.441 \times 10^{-16}$ \\ \hline
        $20$ & 2.7182818284590455 & $4.441 \times 10^{-16}$ \\ \hline
    \end{tabular}
    \vspace{0.2cm}
    \caption{Resultados da aproximação da constante de Euler utilizando a Fórmula B.}
    \label{tab:formula_b}
\end{table}

\end{document}
